Para sintetizar diferentes dimensões de comportamento e participação dos estudantes em métricas quantitativas normalizadas entre 0 e 1, foram criadas as métricas de pontuação. Esses indicadores possibilitam comparações consistentes entre diferentes aulas, turmas e períodos, oferecendo uma visão estruturada da qualidade do engajamento. As métricas apresentadas nesta seção combinam múltiplos aspectos — cognitivos, comportamentais e temporais — por meio de abordagens ponderadas.

\subsection{\textit{Score} Geral de Engajamento}
Essa métrica integra diferentes dimensões de participação em um único valor entre 0 e 1. Busca capturar, de forma sintética, a qualidade do engajamento estudantil combinando indicadores de foco, presença temporal e participação ativa. Essa abordagem ponderada permite uma avaliação holística do comportamento dos alunos durante a aula, fornecendo uma medida única, interpretável e comparável entre sessões distintas. A definição formal é expressa por:
\begin{equation}
\text{Engagement Score} = \sum_{i=1}^{5} w_i \times m_i
\end{equation}

$w_i$ corresponde ao peso atribuído a um componente específico da métrica e cada $m_i$ corresponde ao valor normalizado desse componente. Os pesos são definidos de forma que sua soma seja igual a 1, assegurando que o resultado permaneça dentro do intervalo unitário. Pode-se observar os valores atribuídos a esses pesos na Tabela~\ref{tab:engagement_score_components} a seguir.

\begin{table}[ht]
    \centering
    \vspace{0.5cm}
    \begin{tabular}{|c|c|l|}
        \hline
        Componente & Peso & Descrição \\
        \hline
        Foco & 35\% & Tempo na aba da aula (\texttt{lecture\_focus\_ratio}) \\
        % \hline
        Presença & 30\% & Razão tempo médio/duração total (\texttt{attendance\_ratio}) \\
        % \hline
        Participação & 20\% & Sinais ativos de engajamento (\texttt{voluntary\_participation}) \\
        % \hline
        Câmera & 10\% & Tempo com câmera ativa (\texttt{camera\_engagement}) \\
        % \hline
        Microfone & 5\% & Tempo com microfone ativo (\texttt{mic\_engagement}) \\
        \hline
    \end{tabular}
    \caption{Componentes e pesos do Score Geral de Engajamento}
    \label{tab:engagement_score_components}
\end{table}

A estrutura adotada estabelece um peso maior para fatores diretamente relacionados à atenção contínua e ao tempo efetivamente dedicado à aula, seguidos por indicadores de participação ativa e contribuições audiovisuais. Essa distribuição ponderada privilegia a dimensão cognitiva do engajamento, ao mesmo tempo em que incorpora elementos comportamentais e de interação síncrona.

% \textbf{Justificativa dos pesos}:
% \begin{itemize}
%     \item \textbf{Foco na Aula (35\%)}: Considerado o indicador mais direto de atenção cognitiva
%     \item \textbf{Presença Temporal (30\%)}: Base fundamental para qualquer tipo de engajamento
%     \item \textbf{Participação Voluntária (20\%)}: Indica iniciativa ativa do estudante
%     \item \textbf{Engajamento por Câmera (10\%)}: Contribuição visual limitada por fatores contextuais
%     \item \textbf{Engajamento por Microfone (5\%)}: Participação vocal menos frequente em aulas expositivas
% \end{itemize}

\subsection{Saúde da Atenção}

Avalia a qualidade do padrão de foco dos estudantes ao longo da aula. Diferentemente de medidas puramente quantitativas, esse indicador considera aspectos qualitativos como sustentação do foco, estabilidade temporal e baixa fragmentação, cujo resultado é obtido conforme a Equação~\ref{eq:attention_health}.

\begin{equation}\label{eq:attention_health}
\text{Attention Health} = \sum_{i=1}^{4} w_i \times h_i
\end{equation}

Expresso em um valor entre 0 e 1, \textit{Attention Health} permite comparar a consistência do engajamento entre diferentes aulas e grupos de alunos.

Para cada $w_i$ — peso atribuído a uma dimensão de saúde da atenção — e cada $h_i$ — valor normalizado dessa dimensão — os pesos são definidos de forma que sua soma seja igual a 1, como indicado na Tabela~\ref{tab:attention_health_components}.

\begin{table}[H]
    \centering
    \vspace{0.5cm}
    \begin{tabular}{|c|c|l|}
        \hline
        Dimensão & Peso & Descrição \\
        \hline
        Duração do Foco & 30\% & Capacidade de manter foco sustentado \\
        Baixa Fragmentação & 25\% & Minimização de interrupções no foco \\
        Controle de Multitarefa & 25\% & Gestão adequada de múltiplas tarefas\\
        Consistência Temporal & 20\% & Estabilidade do engajamento ao longo do tempo\\
        \hline
    \end{tabular}
    \caption{Dimensões e pesos da Saúde da Atenção}
    \label{tab:attention_health_components}
\end{table}

Assim, a combinação ponderada das quatro dimensões permite capturar não apenas quanto tempo o aluno permaneceu focado, mas também como esse foco se distribuiu ao longo da sessão, oferecendo uma métrica mais adequada para análises de engajamento.

% \lstinputlisting[
%     language=python,
%     caption={Cálculo da Saúde da Atenção},
%     firstline=49,
%     lastline=110,
%     label={lst:calculate_attention_health}
% ]{scripts/metricsCalc/scores.py}


% \textbf{Funções auxiliares}:
% \begin{itemize}
%     \item \textbf{calculate\_focus\_duration\_health()}: Avalia a adequação da duração dos períodos de foco (nem muito curtos, nem muito longos)
%     \item \textbf{calculate\_fragmentation\_health()}: Mede o grau de continuidade atencional (menor fragmentação = mais saudável)
%     \item \textbf{calculate\_multitasking\_health()}: Avalia a gestão de múltiplas tarefas (multitasking produtivo vs. destrutivo)
%     \item \textbf{calculate\_consistency\_health()}: Analisa a estabilidade do engajamento ao longo da sessão
% \end{itemize}

\subsection{Risco de Distração}
Denota a propensão dos estudantes à desatenção durante as aulas. Este \textit{score}, variando de 0 a 1, integra múltiplos indicadores comportamentais para fornecer uma medida sintética do nível de distração, sendo valores mais altos indicativos de maior propensão à perda de foco. O cálculo descrito na Equação~\ref{eq:distraction_risk} a seguir.

\begin{equation}
\label{eq:distraction_risk}
\text{Distraction Risk} = \sum_{i=1}^{4} w_i \times f_i
\end{equation}

Aqui, $w_i$ representa o peso atribuído ao fator de risco $i$ e $f_i$ corresponde ao valor normalizado desse fator, com pesos definidos de forma que sua soma seja igual a 1. A Tabela~\ref{tab:distraction_risk_components} indica o peso atribuído a cada propriedade.

\begin{table}[ht]
    \centering
    \vspace{0.5cm}
    \begin{tabular}{|c|c|l|}
        \hline
        Fator de Risco & Peso & Descrição \\
        \hline
        Proporção de Distração & 40\% & Percentual de tempo em sites distractores\\
        Frequência de Distração & 25\% & Número de transições para distração por hora\\
        Força das Distrações & 20\% & Impacto dos distractores principais\\
        Intensidade de Trocas & 15\% & Ritmo de alternância entre abas\\
        \hline
    \end{tabular}
    \caption{Fatores de risco e pesos do Score de Distração}
    \label{tab:distraction_risk_components}
\end{table}

A escolha dos pesos reflete a importância relativa de cada elemento observado durante a aula, privilegiando aqueles que apresentam maior correlação com episódios de perda de atenção. 

% \lstinputlisting[
%     language=python,
%     caption={Cálculo do Risco de Distração},
%     firstline=170,
%     lastline=210,
%     label={lst:calculate_distraction_risk}
% ]{scripts/metricsCalc/scores.py}



% \textbf{Funções auxiliares}:
% \begin{itemize}
%     \item \textbf{calculate\_main\_distraction\_strength()}: Calcula o impacto relativo das principais categorias de distração detectadas
%     \item \textbf{calculate\_tab\_switch\_intensity()}: Normaliza e pondera a frequência de troca de abas
% \end{itemize}